\scanpage{315}
\mspage{265}

% slanted
\margin{Réflexions \unsure{provisoires}, \unsure{supplantées} par \uline{XVI}.}

\inserted{(\S 8)} \uline{Appendice} (Retour sur les axiomes $W(7)$, $W(7\ \mathrm{bis})$ (pages 75, 76).

On a vu (prop. 6 p. 78) que $\deleted{\unsure{W(6)}}W(7\ \mathrm{bis})$ implique $W(7)$. Montrons une réciproque partielle :

\uline{Proposition} Supposons $W(7\ \mathrm{bis})$ (axiome \unsure{des|du} \unsure{carré|carrés} \unsure{cocartésien|cocartésiens}) satisfait, et plaçons nous sous les conditions de l'énoncé de $W(7)$, \deleted{\illegible{4}}\inserted{et plus généralement} dans le cas d'un diagramme cubique

\diagram{\(\begin{tikzcd}[ampersand replacement=\&,column sep=small,row sep=small]
{}\&X_2\arrow[rr]\arrow[dd,dashed,"{f_2}" near end]\&{}\&X_3\arrow[dd,"{f_3}"]\\
X_0\arrow[ur,leftrightarrow]\arrow[rr,leftrightarrow]\arrow[dd,"{f_0}"']\&{}\&X_{\unsure{2|1}}\arrow[ur]\arrow[dd,"{f_{\unsure{1|2}}}"']\&{}\\
{}\&Y_2\arrow[rr,dashed]\&{}\&Y_3\\
Y_0\arrow[ur]\arrow[rr]\&{}\&Y_{\unsure{1|2}}\arrow[ur]\&{}
\end{tikzcd},\)}

où les deux carrés $X$ et $Y$ (faces sup. et inf. du cube) sont $W$-cocartésiens. Considérons les conditions suivantes sur ce diagramme

\[
\begin{aligned}
&\text{a)}\ f_0\in W\\
&\text{b)}\ f_1\text{ et }f_2\in W\\
&\text{c)}\ f_3\in W.
\end{aligned}
\]

Alors on a les implications

\[
a\mathbin{\&}b\Longrightarrow c,\qquad b\mathbin{\&}c\Longrightarrow a.
\]

De plus si $W$ est stable par facteurs directs (p. ex. $W$ fortement saturé), alors on a aussi

\[
a\mathbin{\&}c\Longrightarrow b
\]

donc \inserted{la conjonction \unsure{de}} deux \inserted{quelconques} parmi les conditions a) b) c) implique la troisième.

Démonstration. Utilisant la factorisation
% archive sheet number: 315
